\documentclass[11pt]{article}
\usepackage[a4paper,margin=27mm]{geometry}
\usepackage{amsmath,amssymb,mathtools}
\usepackage{fontspec,unicode-math}
\setmainfont{Latin Modern Roman}
\setmathfont{Latin Modern Math}
\usepackage{xurl}
\usepackage[hidelinks]{hyperref}
\newcommand{\sref}[2]{\hyperlink{source:#1:#2}{[S#2]}}
\newcommand{\cataloguescope}[1]{\paragraph{Mathematical question.}#1}
\hypersetup{pdftitle={Problem statement}}
\begin{document}
\section{Agrawal's conjecture}\label{statement}

\subsection{Definitions and mathematical statement}

\paragraph{Shared notation.}$\mathbb N=\{1,2,\ldots\}$, and $\mathbb Z,\mathbb Q,\mathbb R,\mathbb C$ denote the integers, rationals, reals and complex numbers. Any other conventions are specified locally.


For positive coprime integers $n,r$, polynomial congruence modulo $n$ and $X^r-1$ means equality in the quotient ring $(\mathbb Z/n\mathbb Z)[X]/(X^r-1)$; equivalently the difference belongs to the ideal $(n,X^r-1)$ in $\mathbb Z[X]$. Primality always means an integer greater than one with no positive divisors other than one and itself.
\paragraph{Statement recorded in the supplied source.}
Whenever $\gcd(n,r)=1$ and $(X-1)^n\equiv X^n-1\pmod{n,\ X^r-1}$, must either $n$ be prime or $n^2\equiv1\pmod r$? \sref{statement}{2}
\subsection{Short English statement}
Does this single polynomial identity force either primality or the stated square congruence for every coprime pair of positive integers?
\subsection{Sources}
\begin{description}
\item[\hypertarget{source:statement:1}{S1}] ULAM.AI and the UnsolvedMath Contributors, UnsolvedMath: A Curated Collection of Open Mathematics Problems, catalogue entry. CC BY 4.0. \url{https://huggingface.co/datasets/ulamai/UnsolvedMath}.
\item[\hypertarget{source:statement:2}{S2}] Wikipedia contributors, \emph{List of unsolved problems in mathematics}, revision 1366636767 (2026) Number theory, Agrawal's conjecture. \url{https://en.wikipedia.org/w/index.php?title=List_of_unsolved_problems_in_mathematics&oldid=1366636767#Number_theory}.
\item[\hypertarget{source:statement:3}{S3}] Wikipedia contributors, \emph{Agrawal's conjecture}, revision 1295268932 (2026) Formal statement. \url{https://en.wikipedia.org/w/index.php?title=Agrawal%27s+conjecture&oldid=1295268932}.
\end{description}
\label{end:statement}
\end{document}
